Redes varejistas precisam escolher políticas de reposição para demandas heterogêneas. Esta dissertação propõe o \textit{Adaptive Inventory Policy Engine} (AIPE), uma metodologia que organiza cinco etapas: especificar a decisão, representar o contexto, construir evidência operacional comparável, recomendar uma política e avaliar suas consequências. O \textit{Policy Selection Engine} (PSE) é seu componente de escolha, com duas abordagens: associação por perfil informado e estimação por série. A metodologia distingue seleção da família e calibração de parâmetros e admite conservar uma referência quando falta evidência para substituí-la. Sua avaliação utiliza sete políticas em duas aplicações separadas. Na Bahia, 145 séries em duas origens totalizam 4.350 ajustes; a seleção por perfil aumenta o custo médio em 1,93\% e o atendimento em 3,89 pontos percentuais frente à referência global. No piloto M5, 30 séries de dez lojas em três origens totalizam 1.350 ajustes. O preditivo original apresenta maior custo e menor atendimento que a referência. Uma comparação exploratória que reutiliza essas simulações avalia seleção por custo e serviço previstos: com recuo, o custo médio é 1,56\% menor e o atendimento 0,27 ponto maior, mas apenas três decisões mudam e uma série concentra cerca de 92\% da economia. Não se demonstra superioridade generalizável. Oráculos retrospectivos distinguem erros de escolha, insuficiência do portfólio e inviabilidade física. A contribuição é a formulação e implementação do procedimento AIPE, acompanhadas de evidências sobre sua utilidade e seus limites. Coortes retrospectivas, calibrações distintas, horizontes curtos e desenvolvimento no piloto exigem confirmação independente.

\vspace{2em}
\textbf{Palavras-chave}: Metodologia de seleção; Gestão de estoques; Seleção contextual de políticas; Nível de serviço; AIPE.
